\documentclass[12pt]{article}
\usepackage{amsmath,amssymb}

\begin{document}
\section*{{2025 AIME II Problems/Problem 1}}
Source: AoPS Wiki

\subsection*{Solution}
[asy] pair A,B,C,D,E,F,G; A=(0,0); label("$A$", A, S); B=(1.5,0); label("$B$", B, S); C=(2.9,0); label("$C$", C, S); D=(4.2,0); label("$D$", D, S); E=(5.3,0); label("$E$", E, S); F=(6.5,0); label("$F$", F, S); G=(3.7,3); label("$G$", G, N); draw(A--B--C--D--E--F); draw(C--G--D); draw(B--G--E); [/asy] Let $AB=a$ , $BC=b$ , $CD=c$ , $DE=d$ and $EF=e$ . Then we know that $a+b+c+d+e=73$ , $a+b=26$ , $b+c=22$ , $c+d=31$ and $d+e=33$ . From this we can easily deduce $c=14$ and $a+e=34$ thus $b+c+d=39$ . Using Heron's formula we can calculate the area of $\triangle{CGD}$ to be $\sqrt{(42)(28)(12)(2)}=168$ , and since the base of $\triangle{BGE}$ is $\frac{39}{14}$ of that of $\triangle{CGD}$ $\triangle{BGE}$ shares an altitude with $\triangle{CGD}$ , we conclude they are proportional and we can calculate the area of $\triangle{BGE}$ to be $168\times \frac{39}{14}=\boxed{468}$ . ~ Quick Asymptote Fix by eevee9406 , edited by aoum

\end{document}
